\clearpage\thispagestyle{empty}
\pdfbookmark[0]{Quellen zum Ceres-Bericht vom Februar 1802}{ceres180202-sources}
\begin{center}\textbf{Quellen zum Ceres-Bericht vom Februar 1802}\end{center}
\noindent Diese KI-Übertragung in heutiges Deutsch folgt dem deutschen Auszug in Carl Friedrich Gauß’ \href{https://gdz.sub.uni-goettingen.de/id/PPN236007467}{\textit{Werken}, Band~VI} (Göttingen, 1874), S.~204--205. Franz Xaver von Zach zitiert Carl Friedrich Gauß und Heinrich Wilhelm Matthias Olbers in seinem Bericht über die Wiederauffindung der Ceres.
\par\smallskip\noindent Von Zachs \href{https://api.digitale-sammlungen.de/iiif/image/v2/bsb10538597_00178/full/1450,/0/default.jpg}{ursprünglicher Aufsatz}, Nr.~XX, erschien im Februarheft 1802 der Zeitschrift \textit{Monatliche Correspondenz zur Beförderung der Erd- und Himmels-Kunde} (Gotha, Becker), Band~V, S.~170--184. Der Auszug beginnt mit \href{https://api.digitale-sammlungen.de/iiif/image/v2/bsb10538597_00187/full/1450,/0/default.jpg}{Gauß’ Bemerkungen auf S.~179} und endet mit \href{https://api.digitale-sammlungen.de/iiif/image/v2/bsb10538597_00189/full/1450,/0/default.jpg}{Olbers’ Bemerkungen auf S.~181}. Der Aufsatz selbst reicht bis \href{https://api.digitale-sammlungen.de/iiif/image/v2/bsb10538597_00192/full/1450,/0/default.jpg}{S.~184}. Die Digitalisate der Zeitschrift stammen von der Bayerischen Staatsbibliothek.
\par\smallskip\noindent Übertragung und Satz wurden mit KI erstellt.
\clearpage
